%%%%%%%%%%%%%%%%%%%%%%%%%%%%%%%%%%%%%%%%%%%%%%%%%%%%%%%%%%%%%%%
\section{Introduction}\label{secInt}
%%%%%%%%%%%%%%%%%%%%%%%%%%%%%%%%%%%%%%%%%%%%%%%%%%%%%%%%%%%%%%%

Recent work has rapidly advanced the mathematical and logical reasoning abilities of large language models (LLMs) by leveraging \emph{reinforcement learning with verifiable rewards} (RLVR), wherein a policy produces chain-of-thought (CoT) trajectories and receives rule-based, deterministic feedback on final-answer correctness \citep{shao2024deepseekmath,qwen2024qwen25,yu2025dapo,he2025deltal}. A distinctive challenge in RLVR is that response trajectories vary widely in length, often from tens to thousands of tokens and typically increasing as training progresses. So updates computed from a single batch can differ dramatically in both magnitude and variance \citep{yu2025dapo,he2025deltal}. Consequently, training stability and sample efficiency hinge on how we aggregate per-trajectory gradients under this length variability.

A common remedy is \emph{length-dependent} normalization. GRPO normalizes at the sample level by implicitly or explicitly dividing by trajectory length, DAPO normalizes at the batch level by the total length, and Dr.\ GRPO removes length terms to preserve an unbiased estimator \citep{shao2024deepseekmath,yu2025dapo,liu2025drgrpo}. While such strategies can curb instabilities from long CoTs, their statistical properties differ: length-based factors may introduce bias in the policy-gradient direction, whereas length-agnostic updates can suffer high variance—both of which impede convergence \citep{williams1992simple,sutton2018rl,schulman2016gae,he2025deltal}. In particular, estimators with higher gradient variance are more prone to deviating from the optimal optimization direction, resulting in slow or unstable learning.

\paragraph{Related work and connections.}
CAPO is positioned at the intersection of three literatures.

First, it relates to modern RLHF/RLVR optimization pipelines for LLMs, which typically instantiate a policy-gradient method (most often a PPO-type update \citep{schulman2017ppo}) equipped with reward modeling, preference modeling, or verifiable rewards \citep{ouyang2022instructgpt,shao2024deepseekmath,rafailov2023dpo}.
The specific normalization rules compared in this paper ($\Delta L$, GRPO-, DAPO-, and Dr.~GRPO-type variants) can be interpreted as different choices of \emph{trajectory-level reweighting} and \emph{baseline} in a policy-gradient estimator.
Our contribution is to place these heuristics into a single estimation-theoretic framework and to isolate a design axis that is largely implicit in prior work: how \emph{within-trajectory dependence} and heteroscedasticity induced by variable-length CoT should enter the weighting rule.

Second, CAPO has close ties to classical weighted least squares and generalized least squares.
The MVU weights in Lemma~\ref{lem:mvu} are the Gauss--Markov/BLUE solution for estimating a common mean under heteroscedastic noise, and the dependence-aware extensions can be viewed as (structured) GLS with autocorrelated errors \citep{hamilton1994time}.
This connection motivates both the functional form of the optimal weights and the emphasis on correctly modeling how variance grows with length and dependence.

Third, our empirical-Bayes (EB) construction in \S\ref{sec:eb-theory} follows a long tradition of EB modeling for structured variance and shrinkage, in which hyperparameters are fitted by marginal likelihood and then plugged into a decision rule \citep{efron2012large}.
In the present setting, EB serves a pragmatic purpose: it yields closed-form weight expressions with stable large-$L$ behavior while avoiding the computational overhead of full Bayesian integration over dependence parameters.

This paper proposes \emph{Covariance-Aware Policy Optimization (CAPO)}, a principled, variance-aware gradient aggregation framework for RLVR. We begin by casting per-step loss/gradient aggregation as a classical \emph{minimum-variance unbiased estimation} (MVUE) problem over heterogeneous trajectory lengths, in the same spirit as \citet{he2025deltal}. We then introduce a generalized Bayesian working model for trajectory-level gradient statistics, place a Jeffreys prior on the scale parameter to obtain an invariant baseline, and show that the resulting posterior-precision aggregator recovers $\Delta L$ normalization as a special case \citep{jeffreys1946invariant,he2025deltal}. Finally, we extend beyond token-independence by modeling realistic within-trajectory dependencies via simple covariance structures (e.g., banded or geometrically decaying kernels). This yields a plug-in weighting rule that down-weights trajectories by both an empirically estimated variance and a sublinear length factor:
\begin{equation}
\label{eq:capo-agg}
\bar g = \sum_{i=1}^{M} \tilde w_i g_i,
\qquad
\tilde w_i \propto \frac{L_i^{-\alpha}}{\widehat{v}_i},
\qquad
\alpha \in (0,1),
\end{equation}
where $g_i$ is the unnormalized gradient contribution from trajectory $i$, $L_i$ is its length, and $\widehat{v}_i$ estimates the trajectory-level gradient variance under the learned covariance model.

Our analysis shows that, under broad dependence classes, the trajectory-level variance grows \emph{approximately linearly} with length whenever within-trajectory correlations are summable (short-range dependence). For a simple but illustrative case with geometrically decaying correlations $\mathrm{Cov}(Z_t,Z_s)=\sigma^2\rho^{|t-s|}$ (with $|\rho|<1$), one obtains
\begin{equation}
\label{eq:linear-variance}
\mathrm{Var}\left(\sum_{t=1}^{L} Z_t\right)
=
\sigma^2 \sum_{k=-(L-1)}^{L-1} (L-|k|)\,\rho^{|k|}
= \Theta(L),
\qquad |\rho|<1,
\end{equation}
so that naive length-proportional weighting can be unstable. In finite samples, additionally, variance estimation noise increases with $L$, which motivates the sublinear dampening in \eqref{eq:capo-agg}. We further study the \emph{coefficient of variation} (CV) of the aggregated estimator under different normalization schemes and show that CAPO reduces CV relative to length-agnostic updates while avoiding the length-induced bias of sample- or batch-normalization.

Empirically, we evaluate CAPO on a focused RLVR testbed: a representative task (CountDown), a single open-weight backbone model (Qwen2.5--1.5B--Instruct), two context-length regimes (2048 and 8192), and a compact set of aggregation baselines that isolate the effect of aggregation under matched tokens and identical decoding/clipping (GRPO Norm, DAPO Norm, Dr.~GRPO Norm, $\Delta L$, and CAPO). In addition to end performance, we report \emph{stability diagnostics} (monotonicity of held-out accuracy, step-to-step drift, CV/kurtosis of trajectory weights, and length-decile stratification), together with a targeted ablation over the covariance bandwidth used by CAPO's empirical-Bayes estimator.

\paragraph{Summary of contributions.}
(i) A unified estimation-theoretic and Bayesian analysis of RLVR gradient aggregation that derives $\Delta L$ from MVUE/Jeffreys principles and clarifies invariances and bias--variance trade-offs \citep{jeffreys1946invariant,he2025deltal}; (ii) \emph{CAPO}, a covariance-aware generalization with consistent variance estimation from rollouts and closed-form plug-in weights as in \eqref{eq:capo-agg}; (iii) dependence-aware theory linking token-level autocovariance to trajectory-level variance growth, yielding principled corrections beyond length-only normalization; and (iv) an empirical protocol with stability diagnostics and targeted ablations that isolate covariance-aware weighting effects under matched compute.

\paragraph{Organization.}
Section~\ref{sec:setup} formalizes the trajectory aggregation problem and clarifies the role of the scaling constraint in \eqref{eq:linear-agg}.
Section~\ref{sec:method} develops the MVU and Bayesian foundations, derives dependence-aware variance laws, and establishes the EB estimator and its asymptotics.
Section~\ref{sec:algorithms} provides implementable CAPO variants and complexity considerations.
Section~\ref{sec:results} describes the RLVR evaluation protocol and stability diagnostics.
Appendices supply full proofs and illustrative covariance examples.
